
\section{Introduction}

Coding agents now work at repository scale. They inspect code, edit files, run tests, collect logs, and synthesize patches across many steps. That workflow exposes a recurring weakness: each session begins with partial amnesia. The next agent may inherit the codebase, but not the local truth that made the last session coherent.


\begin{figure}[H]
\centering
\begin{tikzpicture}[
  every node/.style={figlabel},
  panel/.style={figbox, minimum width=5.35cm, minimum height=5.05cm, align=center},
  loop/.style={draw=BitSlate!85, line width=0.85pt, circle, minimum size=1.2cm, align=center}
]
\node[panel] (left) at (0,0) {};
\node[anchor=north, figtitle, font=\sffamily\bfseries\large] at ($(left.north)+(0,-0.26)$) {Without CONTINUITY};
\foreach \y/\txt in {1.20/{observe},0.0/{reason},-1.20/{act}} {
  \node[loop] at (-0.95,\y) {\txt};
}
\draw[BitRose!90!black, line width=1.6pt] ($(left.west)+(0.62,0.58)$) -- ($(left.east)+(-0.62,0.58)$);
\node[figsmall, text=BitRose!90!black, fill=white, inner sep=1pt] at (0.0,0.86) {session boundary};
\draw[figarrow] (-0.95,0.6) -- (-0.95,0.32);
\draw[figarrow] (-0.95,-0.6) -- (-0.95,-0.88);
\node[align=center, figsmall] at (0,-2.02) {work restarts\\context evaporates};

\node[panel, fill=BitBlue!6, draw=BitBlue!65] (right) at (7.25,0) {};
\node[anchor=north, figtitle, font=\sffamily\bfseries\large] at ($(right.north)+(0,-0.26)$) {With CONTINUITY};
\draw[BitBlue!78, line width=1.8pt] ($(right.center)+(-1.65,-1.65)$)
   .. controls +(0,1.05) and +(0,-1.0) .. ($(right.center)+(-0.75,-0.95)$)
   .. controls +(0,1.0) and +(0,-1.0) .. ($(right.center)+(0.0,-0.18)$)
   .. controls +(0,1.0) and +(0,-1.0) .. ($(right.center)+(0.78,0.55)$)
   .. controls +(0,1.0) and +(0,-1.0) .. ($(right.center)+(1.58,1.6)$);
\foreach \x/\y in {-1.65/-1.65,-0.75/-0.95,0/-0.18,0.78/0.55,1.58/1.60} {
  \fill[white] ($(right.center)+(\x,\y)$) circle (0.13cm);
  \draw[BitBlue!85, line width=0.75pt] ($(right.center)+(\x,\y)$) circle (0.13cm);
}
\node[livebox, minimum width=1.65cm, minimum height=.52cm, font=\sffamily\scriptsize\bfseries] at ($(right.center)+(-1.7,1.52)$) {\statefile{STATE.md}};
\node[align=center, figsmall] at ($(right.center)+(0,-2.03)$) {work compounds\\the loop resumes};

\draw[figflow, line width=1.35pt] (left.east) -- node[above, figsmall, fill=white, inner sep=1.5pt] {persist state} (right.west);
\end{tikzpicture}
\caption{A session can end without erasing the work state. That is the core operational promise.}
\label{fig:intro-state-gap}
\end{figure}


Four kinds of state are commonly lost. First, \term{current truth} drifts: the repository may have changed, but the agent still acts on an earlier plan. Second, \term{attempt history} disappears: failures are repeated because earlier experiments were never captured. Third, \term{operational constraints} get buried in chat logs, issue threads, or a developer's head. Fourth, \term{handoff state} is weak: one agent cannot easily tell another what remains, what is blocked, and what evidence already exists.

Tool practice has started to converge on a partial fix. Codex loads \statefile{AGENTS.md} before work and composes instructions across global, repository, and directory scopes \cite{OpenAICodexAgents}. Claude Code documents persistent repository context through \file{CLAUDE.md} and distinguishes that soft context from hard enforcement \cite{AnthropicClaudeMemory,AnthropicClaudeHooks}. Hermes progressively discovers local context files as the agent moves through a project \cite{HermesContextFiles}. OpenClaw stores memory as plain Markdown on disk rather than hidden platform state \cite{OpenClawMemory}. These systems already treat the repository as part of the agent interface.

\systemname{} extends that interface into a compact state substrate. It does not ask one file to do everything. It separates stable boot guidance from mutable state, active tasks, learned lessons, append-only history, and generated evidence. The minimal core stays intentionally small:

\begin{calloutbox}[colback=BitGray,colframe=BitLine,title=Minimal CONTINUITY core]
\begin{tabularx}{\linewidth}{>{\raggedright\arraybackslash}p{2.9cm}X}
\statefile{AGENTS.md} & boot order, governance, load rules \\
\statefile{STATE.md} & current project truth \\
\statefile{TODO.md} & executable task queue \\
\statefile{LEARNINGS.md} & compact reflection and reusable lessons \\
\statefile{MEMORY.md} & durable facts index, loaded selectively \\
\file{completed/} & append-only records of finished work \\
\file{snapshots/} & recovery checkpoints \\
\file{artifacts/} & outputs, logs, test evidence \\
\statefile{CHECKPOINT.md} & approval gates for high-risk actions \\
\end{tabularx}
\end{calloutbox}

The point is not to create a repository diary. The point is to keep the smallest set of files that preserve continuity across sessions and across agents.

Recent evidence makes the design constraint sharper. Gloaguen et al. find that repository context files do not reliably improve coding-agent task success and can increase inference cost by more than 20\% on average \cite{Gloaguen2026Agents}. Dos Santos et al. report six recurring configuration smells in \statefile{AGENTS.md}-style files, including \term{context bloat} (too much irrelevant or diffuse instruction), \term{lint leakage} (tool-output noise or rules pasted directly into the manifest), \term{skill leakage} (specialized or personal instructions spilling into shared context), and \term{conflicting instructions} \cite{DosSantos2026Smells}. Those results rule out the obvious mistake: more prose is not automatically better state.

The rest of the paper develops the positive version of that negative lesson. \systemname{} works when it remains compact, scoped, current, and reviewable. The next sections show where that stance comes from, how the file grammar works, how multi-agent collaboration fits, and how the design should be evaluated and constrained.
